\section{Conclusion}

The paper asks what changes when current AI capability augments not only
inference compute but also compute used to improve AI. The answer is not a
notational detail. It changes the relative price of automated research services,
adds a capability term to the developer's costate equation, and strengthens the
feedback from AI improvement to future AI improvement.

Three results organize the analysis. First, the research elasticity determines
whether a rising wage relative to the price of automated research services makes
automated research asymptotically dominant.
Second, the final-production elasticity determines whether human labor remains a
long-run bottleneck. With complementarity, output per person cannot keep growing
along the balanced path characterized here when production labor remains a
positive population share. At unit
elasticity, a balanced-growth candidate with positive capability growth can be
feasible only when decreasing returns to effective research dominate the
combined effect of capability on inference and automated research. With gross
substitution, if AI services asymptotically dominate final production, unbounded
capability makes \(Y/K\) and the return to capital
unbounded, ruling out finite constant growth rates.

Third, factor shares and factor prices must be distinguished. Along the
AI-dominated accelerating limit, labor's income share tends to zero. Yet the real
wage eventually rises when final-production substitution is below
\(1/(\alpha\eta)\) along the characterized limit; above that threshold the
conditional wage result reverses. A vanishing share does not by itself imply a
falling level when the scale of the economy grows sufficiently fast.
At the threshold itself, the leading asymptotic term is zero and the model does
not sign the eventual wage change.

The perfect-foresight simulations reinforce the analytical results without being
used as proofs. At unit elasticity in final production, both Cobb--Douglas and
gross-substitutes research paths converge to the rates implied by the model.
When automated research services are gross substitutes for human research, the
human-research employment share is hump-shaped: it can increase for centuries in
the numerical normalization before automation eventually dominates. Separating
\(M\) from \(AM\) makes this transition visible. The simulations also confirm
the Euler-equation link between consumption growth and the net real interest
rate. These paths are numerical illustrations, not forecasts of calendar time.

Several extensions are important. A direct standing-on-shoulders term could be
added to determine how much it amplifies the feedback already carried by \(AM\).
Physical capital could enter AI research explicitly, separating chips,
datacenters, energy infrastructure, and model capability. Alternative market
structures could clarify how competition changes the balance between current
inference supply and research investment.

Finally, greater automation of AI research may create agency and supervision
problems that are absent here. If human oversight is needed to keep automated
research aligned with the developer's objective, an integrated monopolist may
internalize part of that problem and moderate research speed. Competing developers
may instead reduce supervision to avoid losing a race. That extension would make
human research valuable not only as a productive input but also as a governance
input. It could provide a disciplined basis for studying subsidies to human AI
research or taxes on insufficiently supervised research automation; those policy
claims require an explicit agency model and do not follow from the present
technology alone.
